\scanpage{236}
\mspage{$120'$}

2. Condition d'isomorphie pour une donnée $\pi$.

\deleted{On introduit}

\[
\deleted{(45)\qquad G=\{g\in A=\mathrm{Aut}(\Gamma)\mid \substack{g\text{ centralise }Z\ \text{i.e. }g\in\mathcal G\\g_{\unsure{\mathcal G}}(\pi)=}\,}
\]

On s'intéresse aux automorphismes d'une structure \deleted{$\underset{\substack{\shortparallel\\(S,\check S,R)}}{\Gamma},Z\subset\unsure{G}_{S}\ \unsure{Q}_{\check S}$}

\[
M=\bigl(\underset{\substack{\shortparallel\\(S,\check S,R}}{\Gamma},\quad Z\subset\mathcal A=\mathrm{Aut}(\Gamma),\quad\pi\subset\Pi_Z\subset\mathrm{Bij}(S,\check S)\bigr),
\tag*{(49)}
\]

% slanted
\margin{\textquotedblleft{}\unsure{prégraphisme|prégraphe} \deleted{\unsure{groupe}} multijaugé\textquotedblright}

ce sont donc les $g\in\mathcal A$ qui \uline{normalisent} $Z$, et qui satisfont

\[
\underbrace{g(\pi)}_{\text{tr. de str.}}\left(\overset{\mathrm{def}}{=}\underbrace{\{g_{\check S}\circ\varphi\circ g_S^{-1}\mid\varphi\in\pi\}}\right)=\pi
\tag*{(50)}
\]

ce qui équivaut à la condition, \deleted{\unsure{que}} \inserted{en termes d'\uline{un} $\varphi_0\in\pi$}

\[
g(\varphi_0)\left(\overset{\mathrm{def}}{=}g_{\check S}\varphi_0g_{\check S}\varphi^{-1}\right)\in\pi\quad.
\tag*{(51)}
\]

On a donc des hom de groupes

\diagram{\begin{tikzcd}[ampersand replacement=\&,column sep=large]
{(52)}\&\mathrm{Aut}(M)\arrow[r]\&\mathrm{Aff}_{Z}(\pi)\arrow[r]\&\mathrm{Aut}(Z),\\
\&\mathcal A\arrow[u,hook]\&\&
\end{tikzcd}}

et je m'intéresse aux autom. $g$ de $M$ dont l'image dans $\mathrm{Aut}(Z)$ soit égale à $1$, i.e. \inserted{qui \uline{centralisent} $Z$} i.e. qui opèrent sur $\pi$ par translation, \deleted{On} donc on s'intéresse au groupe formé de ces $g$ :

\[
G=\left\{g\in\mathcal G\overset{\mathrm{def}}{=}\mathrm{Centr}_{\mathcal A}(Z)\ \middle|\ \substack{g(\pi)=\pi\quad\mathrm{i.e.}\\g_{\unsure{S|\check S}}\varphi_0g_S^{-1}=\varphi_0}\right\}.
\tag*{(53)}
\]

On trouve donc un diagramme canonique.

% archive sheet number: 235

\mspage{119 bis}
\unsure{graphes|graphe}, \inserted{avec} groupe d'homothéties $Z$ \unsure{qu'on se le donne|qu'on se donne|opérant librement}, \inserted{et satisfaisant la condition (41\unsure{*|''})}, \deleted{revient à celle}
